\section{Численный метод решения задачи}
\label{sec:Chapter2} \index{Chapter2}
Содержание данного пункта основано на материале книги \cite{Chymi}.

Для численного решения задачи введем на $\Omega$ сетку $\omega_{h\tau} = \overline{\omega}_h \times \omega_\tau$, где

\[
    T = T_0,
\]
\[
    L_x = L_{x0}, L_y = L_{y0}, L_z = L_{z0}
\]
\[
    \overline{\omega}_h = \{\left(x_i=ih_x, y_j=jh_y, z_k=kh_z\right), i,j,k=0,1,\dots,N, h_xN=L_x,h_yN=L_y,h_zN=L_z\},
\]
\[
    \omega_\tau = \{t_n=n\tau,n=0,1,\dots,K,\tau K=T\}.
\]
Через $\omega_{h}$ обозначим множество внутренних, а через $\gamma_{h}$ множество граничных узлов сетки $\overline{\omega}_{h}$.

Для аппроксимации исходного уравнения (\ref{equation:1}) c однородными граничными условиями (\ref{equation:4})--(\ref{equation:6}) и начальными условиями (\ref{equation:2})--(\ref{equation:3}) воспользуемся следующей системой уравнений:
\[
    \frac{
        u^{n+1}_{ijk}
        -
        2 u^{n}_{ijk}
        +
        u^{n-1}_{ijk}
    }
    {
        \tau^2
    }
    =
    \Delta_h u^{n},
    \qquad
    (x_i,y_j,z_k)\in \omega_h,
    \qquad
    n = 1,2,\dots,K-1,
\]
Здесь \(\Delta_h\) — семиточечный разностный аналог оператора Лапласа:
\[
\Delta_h u^{n}
=
\frac{
        u^{n}_{i-1,j,k}
        -
        2u^{n}_{i,j,k}
        +
        u^{n}_{i+1,j,k}
    }
    {h^2}
+
\frac{
        u^{n}_{i,j-1,k}
        -
        2u^{n}_{i,j,k}
        +
        u^{n}_{i,j+1,k}
    }
    {h^2}
+
\frac{
        u^{n}_{i,j,k-1}
        -
        2u^{n}_{i,j,k}
        +
        u^{n}_{i,j,k+1}
    }
    {h^2}.
\]
Приведенная выше разностная схема является явной -- значения \(u^{n+1}_{ijk}\) на \((n+1)\)-м шаге можно явным образом выразить через значения на предыдущих слоях.

Для начала счета (т.е. для нахождения \(u^{2}_{ijk}\)) должны быть заданы значения \(u^{0}_{ijk}\), \(u^{1}_{ijk}\), \((x_i,y_j,z_k) \in \omega_{h}\).
Из условия (\ref{equation:2}) имеем
\begin{align}
    u^{0}_{ijk}
    =
    \varphi(x_i,y_j,z_k),
    \qquad
    (x_i,y_j,z_k) \in \omega_{h}.
    \label{equation:10}
\end{align}
Простейшая замена начального условия (\ref{equation:3}) уравнением \((u^{1}_{ijk}-u^{0}_{ijk})/\tau=0\) имеет лишь первый порядок аппроксимации по \(\tau\).
Аппроксимацию второго порядка по \(\tau\) и \(h\) дает разностное уравнение
\begin{align}
    \dfrac{u^{1}_{ijk}-u^{0}_{ijk}}{\tau}
    =
    \dfrac{\tau}{2}\Delta_h\varphi(x_i,y_j,z_k),
    \qquad
    (x_i,y_j,z_k) \in \omega_{h}.
    \label{equation:11}
    \\
    u^{1}_{ijk}
    =
    u^{0}_{ijk}+\dfrac{\tau^2}{2}\Delta_h\varphi(x_i,y_j,z_k)
    \label{equation:12}
\end{align}

Разностная аппроксимация для периодических граничных условий выглядит следующим образом
\[
    u^{n+1}_{0jk}=u^{n+1}_{Njk},
    \qquad
    u^{n+1}_{1jk}=u^{n+1}_{N+1jk},
\]
\[
    u^{n+1}_{i0k}=u^{n+1}_{iNk},
    \qquad
    u^{n+1}_{i1k}=u^{n+1}_{iN+1k},
\]
\[
    u^{n+1}_{ij0}=u^{n+1}_{ijN},
    \qquad
    u^{n+1}_{ij1}=u^{n+1}_{ijN+1},
\]
\(i,j,k=0,1,\dots,N\).

Для вычисления значений \(u^0, u^1, \in \gamma_{h}\) допускается использование аналитического значения \(u\), которое задается в программе еще для вычисления погрешности решения задачи.
